\subsection{Whole Gauss coefficient interchange and genuine projective model logs}
\label{subsec:polynomial-coefficient-gauss-models}

\begin{theorem}[True coefficient/variable interchange on whole polynomials]
\label{polynomial:coefficient-variable-interchange}
\lean{EnvelopeProperty.polynomialCoefficientSwap,
EnvelopeProperty.polynomialCoefficientSwap_X,
EnvelopeProperty.polynomialCoefficientSwap_C,
EnvelopeProperty.polynomialCoefficientSwap_monomial,
EnvelopeProperty.polynomialCoefficientSwap_monomial_coeff,
EnvelopeProperty.polynomialCoefficientSwap_coeff_homogeneousComponent}
\leanok
For arbitrary commutative coefficient semirings there is an actual algebra
isomorphism between multivariable polynomials with polynomial coefficients
and polynomials with multivariable coefficients. It sends each entire
coefficient polynomial to its true coefficientwise constant-polynomial
image and each multivariable variable to its constant polynomial image.
Its univariate coefficients commute with every original multivariable
homogeneous-component operator.
\end{theorem}
\begin{proof}
Compose the actual Mathlib polynomial/optional-variable algebra isomorphisms
in their appropriate inverse directions. These actual isomorphisms have
whole-ring inverse laws. Their true constant and variable formulas imply
the asserted coefficient-polynomial and variable identities by polynomial
algebra-homomorphism extensionality and actual coefficient commutation.
An entire multivariable monomial with coefficient polynomial $q$ maps to
the true coefficientwise image of $q$, multiplied by the constant
multivariable monomial. Taking its univariate coefficient gives the
original multivariable monomial with the corresponding actual coefficient
of $q$. A homogeneous-component operator retains or kills that entire
monomial according to its original multivariable degree. Thus both
orders of the claimed operations agree on every genuine monomial;
additive polynomial induction proves agreement on every whole polynomial.
No generator-value norm inference or ring-interchange certificate is used.
\end{proof}

\begin{theorem}[Generic whole Gauss naturality and grading stability]
\label{polynomial:gauss-coefficient-grading-stability}
\lean{EnvelopeProperty.Berkovich.Spectrum.gauss_comap,
EnvelopeProperty.Berkovich.Spectrum.polynomialGauss,
EnvelopeProperty.Berkovich.Spectrum.polynomialGauss_val_X,
EnvelopeProperty.Berkovich.Spectrum.polynomialGauss_val_C,
EnvelopeProperty.Berkovich.Spectrum.polynomialGauss_isPolynomialGraded,
EnvelopeProperty.Berkovich.Spectrum.continuous_polynomialGauss_prod,
EnvelopeProperty.Berkovich.Spectrum.isBounded_gauss,
EnvelopeProperty.Berkovich.NormalizedProjectiveSpectrum.polynomialGauss,
EnvelopeProperty.Berkovich.NormalizedProjectiveSpectrum.continuous_polynomialGauss_prod,
EnvelopeProperty.Berkovich.RelativeNormalizedProjectiveSpectrum.polynomialGauss,
EnvelopeProperty.Berkovich.RelativeNormalizedProjectiveSpectrum.polynomialGauss_val_apply,
EnvelopeProperty.Berkovich.RelativeNormalizedProjectiveSpectrum.continuous_polynomialGauss}
\leanok
\uses{polynomial:coefficient-variable-interchange,spectrum:gauss-point,spectrum:gauss-joint-continuity,projective:general-normalized-spectrum,projective:general-relative-normalized-charts}
An actual ring homomorphism commutes with genuine Gauss pullback for
arbitrary seminorms and nonnegative radii, including radius zero.
Interchanging coefficients in the Gauss extension of a genuinely graded
multivariable seminorm produces another genuinely graded multivariable
seminorm. It retains all original variable values, and its entire
coefficient restriction is the genuine Gauss extension of the original
coefficient seminorm. Normalization therefore gives a true relative
normalized projective point for that precise Gauss coefficient seminorm.
\end{theorem}
\begin{proof}
The two Gauss restrictions have identical original weighted coefficient
values, by the actual polynomial coefficient-map formula. The genuine
Gauss coefficient-bound criterion and lower bound for every coefficient
prove both whole seminorm inequalities, identifying the entire points
without assuming coefficient-map injectivity or finite supports survive.

For the interchanged Gauss point, the preceding entire coefficient and
homogeneous-component commutation makes each projected weighted coefficient
value at most its original weighted coefficient value, by genuine
original gradedness. Each is bounded by the full original Gauss norm.
The genuine Gauss bound criterion then bounds the entire homogeneous
projection by the entire point norm. The independently proved projection
criterion yields true gradedness on every whole polynomial. Constant
variable values follow from the actual algebra interchange and Gauss
constant formula. The entire coefficient-polynomial formula and the
preceding Gauss naturality identify the prescribed coefficient
restriction; normalized variables give the asserted actual projective
point. No field, boundedness, completeness or positive-radius hypothesis
is needed for these generic identities. Whole-ring pullback is continuous
in the pointwise topology, and the genuine Gauss construction is jointly
continuous in the point and radius. Restricting these actual continuous
maps to the derived normalized and fixed-coefficient subtypes proves
the stated continuity. Finally, if every original coefficient value
and the radius are at most one, every weighted coefficient value is
at most one. The true Gauss bound criterion proves whole-point boundedness.
\end{proof}

\begin{theorem}[Actual bounded polynomial coefficient model and original Gauss ideal function]
\label{projective:gauss-polynomial-model-ideal-log}
\lean{EnvelopeProperty.ProjBerk.gaussModelPoint,
EnvelopeProperty.ProjBerk.gaussModelPoint_val_apply,
EnvelopeProperty.ProjBerk.gaussModelCoefficient_le_one,
EnvelopeProperty.ProjBerk.continuous_gaussModelPoint,
EnvelopeProperty.ProjBerk.phiNorm_eq_nativeModelIdealNorm,
EnvelopeProperty.ProjBerk.phiIdeal_eq_nativeModelIdealLog}
\leanok
\uses{polynomial:gauss-coefficient-grading-stability,projective:native-homogeneous-model-ideal-representation,projective:ideal-spectrum-bridge,projective:paper-normalized-spectrum-bridge,lem:general-ideal-pullback}
For the old normalized projective point over a trivially valued field,
the actual Gauss extension of radius $e^{-1}$ defines, after the true
coefficient interchange, a native normalized projective point over the
MODEL coefficient ring $k[t]$ with its genuine bounded Gauss seminorm.
For an actual homogeneous ideal $J\subseteq k[t][X_0,\ldots,X_n]$,
let $I$ be its entire original image in $(k[X_0,\ldots,X_n])[t]$ under
the preceding true interchange. The old whole Gauss ideal norm and
$\phi_I$ equal the intrinsic actual proper-center model norm and logarithm
of the genuine projective kernel ideal sheaf attached to $J$.
\end{theorem}
\begin{proof}
The old point specializes the genuine trivial-coefficient normalized
point, with every original polynomial value retained. The preceding
Gauss stability and entire coefficient identity give the true relative
normalized point over the actual polynomial coefficient ring $k[t]$.
The original coefficient values on $k$ are at most one and $e^{-1}\le1$;
the genuine coefficientwise Gauss bound criterion therefore bounds every
actual coefficient polynomial in $k[t]$. This is a bound on the model
ring, not on a nontrivially valued coefficient field.

Apply the proved actual center-stalk homogeneous model-ideal representation
to this genuine normalized projective point and original ideal $J$.
The actual ideal-map identity under the entire ring interchange, evaluated
at the proved bounded old Gauss point, identifies its whole original
ideal norm with that of $I$. The existing old ideal-bridge producer
identifies this with the original old Gauss ideal norm. Taking the true
extended logarithm and the old proved log-of-sup identity gives $\phi_I$,
including zero ideals and negative infinity. Every equality concerns
entire original polynomial evaluations and whole ideal suprema; original
powers and nonradical ideals are retained.

This identifies the actual projective model over the genuine coefficient
ring and its point/ideal functions. It does not assert a scheme product
isomorphism with $\mathbb P^n_k\times\mathbb A^1_k$ or a coherent analytic
structure-sheaf comparison before constructing those separate producers.
\end{proof}
